% Options for packages loaded elsewhere
% Options for packages loaded elsewhere
\PassOptionsToPackage{unicode}{hyperref}
\PassOptionsToPackage{hyphens}{url}
%
\documentclass[
  letterpaper,
]{book}
\usepackage{xcolor}
\usepackage{amsmath,amssymb}
\setcounter{secnumdepth}{5}
\usepackage{iftex}
\ifPDFTeX
  \usepackage[T1]{fontenc}
  \usepackage[utf8]{inputenc}
  \usepackage{textcomp} % provide euro and other symbols
\else % if luatex or xetex
  \usepackage{unicode-math} % this also loads fontspec
  \defaultfontfeatures{Scale=MatchLowercase}
  \defaultfontfeatures[\rmfamily]{Ligatures=TeX,Scale=1}
\fi
\usepackage{lmodern}
\ifPDFTeX\else
  % xetex/luatex font selection
\fi
% Use upquote if available, for straight quotes in verbatim environments
\IfFileExists{upquote.sty}{\usepackage{upquote}}{}
\IfFileExists{microtype.sty}{% use microtype if available
  \usepackage[]{microtype}
  \UseMicrotypeSet[protrusion]{basicmath} % disable protrusion for tt fonts
}{}
\makeatletter
\@ifundefined{KOMAClassName}{% if non-KOMA class
  \IfFileExists{parskip.sty}{%
    \usepackage{parskip}
  }{% else
    \setlength{\parindent}{0pt}
    \setlength{\parskip}{6pt plus 2pt minus 1pt}}
}{% if KOMA class
  \KOMAoptions{parskip=half}}
\makeatother
% Make \paragraph and \subparagraph free-standing
\makeatletter
\ifx\paragraph\undefined\else
  \let\oldparagraph\paragraph
  \renewcommand{\paragraph}{
    \@ifstar
      \xxxParagraphStar
      \xxxParagraphNoStar
  }
  \newcommand{\xxxParagraphStar}[1]{\oldparagraph*{#1}\mbox{}}
  \newcommand{\xxxParagraphNoStar}[1]{\oldparagraph{#1}\mbox{}}
\fi
\ifx\subparagraph\undefined\else
  \let\oldsubparagraph\subparagraph
  \renewcommand{\subparagraph}{
    \@ifstar
      \xxxSubParagraphStar
      \xxxSubParagraphNoStar
  }
  \newcommand{\xxxSubParagraphStar}[1]{\oldsubparagraph*{#1}\mbox{}}
  \newcommand{\xxxSubParagraphNoStar}[1]{\oldsubparagraph{#1}\mbox{}}
\fi
\makeatother


\usepackage{longtable,booktabs,array}
\newcounter{none} % for unnumbered tables
\usepackage{calc} % for calculating minipage widths
% Correct order of tables after \paragraph or \subparagraph
\usepackage{etoolbox}
\makeatletter
\patchcmd\longtable{\par}{\if@noskipsec\mbox{}\fi\par}{}{}
\makeatother
% Allow footnotes in longtable head/foot
\IfFileExists{footnotehyper.sty}{\usepackage{footnotehyper}}{\usepackage{footnote}}
\makesavenoteenv{longtable}
\usepackage{graphicx}
\makeatletter
\newsavebox\pandoc@box
\newcommand*\pandocbounded[1]{% scales image to fit in text height/width
  \sbox\pandoc@box{#1}%
  \Gscale@div\@tempa{\textheight}{\dimexpr\ht\pandoc@box+\dp\pandoc@box\relax}%
  \Gscale@div\@tempb{\linewidth}{\wd\pandoc@box}%
  \ifdim\@tempb\p@<\@tempa\p@\let\@tempa\@tempb\fi% select the smaller of both
  \ifdim\@tempa\p@<\p@\scalebox{\@tempa}{\usebox\pandoc@box}%
  \else\usebox{\pandoc@box}%
  \fi%
}
% Set default figure placement to htbp
\def\fps@figure{htbp}
\makeatother





\setlength{\emergencystretch}{3em} % prevent overfull lines

\providecommand{\tightlist}{%
  \setlength{\itemsep}{0pt}\setlength{\parskip}{0pt}}



 


\makeatletter
\@ifpackageloaded{bookmark}{}{\usepackage{bookmark}}
\makeatother
\makeatletter
\@ifpackageloaded{caption}{}{\usepackage{caption}}
\AtBeginDocument{%
\ifdefined\contentsname
  \renewcommand*\contentsname{Table of contents}
\else
  \newcommand\contentsname{Table of contents}
\fi
\ifdefined\listfigurename
  \renewcommand*\listfigurename{List of Figures}
\else
  \newcommand\listfigurename{List of Figures}
\fi
\ifdefined\listtablename
  \renewcommand*\listtablename{List of Tables}
\else
  \newcommand\listtablename{List of Tables}
\fi
\ifdefined\figurename
  \renewcommand*\figurename{Figure}
\else
  \newcommand\figurename{Figure}
\fi
\ifdefined\tablename
  \renewcommand*\tablename{Table}
\else
  \newcommand\tablename{Table}
\fi
}
\@ifpackageloaded{float}{}{\usepackage{float}}
\floatstyle{ruled}
\@ifundefined{c@chapter}{\newfloat{codelisting}{h}{lop}}{\newfloat{codelisting}{h}{lop}[chapter]}
\floatname{codelisting}{Listing}
\newcommand*\listoflistings{\listof{codelisting}{List of Listings}}
\makeatother
\makeatletter
\makeatother
\makeatletter
\@ifpackageloaded{caption}{}{\usepackage{caption}}
\@ifpackageloaded{subcaption}{}{\usepackage{subcaption}}
\makeatother
\usepackage{bookmark}
\IfFileExists{xurl.sty}{\usepackage{xurl}}{} % add URL line breaks if available
\urlstyle{same}
% fallback for those not using the hyperref driver hyperxmp:
\makeatletter
\@ifundefined{xmpquote}{\newcommand{\xmpquote}[1]{#1}}{}
\makeatother
\hypersetup{
  pdftitle={NURS 180: Stress and Strategies to Promote Well Being},
  pdfauthor={Corneille Ntihabose},
  hidelinks,
  pdfcreator={LaTeX via pandoc}}


\title{NURS 180: Stress and Strategies to Promote Well Being}
\usepackage{etoolbox}
\makeatletter
\providecommand{\subtitle}[1]{% add subtitle to \maketitle
  \apptocmd{\@title}{\par {\large #1 \par}}{}{}
}
\makeatother
\subtitle{A Student Guide to Stress, Self-Care, and Resilience}
\author{Corneille Ntihabose}
\date{}
\begin{document}
\frontmatter
\maketitle

\renewcommand*\contentsname{Table of contents}
{
\setcounter{tocdepth}{2}
\tableofcontents
}

\mainmatter
\bookmarksetup{startatroot}

\chapter{Welcome to NURS 180}\label{welcome-to-nurs-180}

\bookmarksetup{startatroot}

\chapter{Stress and Strategies to Promote Well
Being}\label{stress-and-strategies-to-promote-well-being}

Welcome to \emph{NURS 180: Stress and Strategies to Promote Well Being}.

This online book introduces foundational concepts and practical
strategies for understanding stress, strengthening resilience, and
promoting physical, mental, emotional, and spiritual well-being.

\section{About this book}\label{about-this-book}

The book follows the course topics outlined in the NURS 180 syllabus. It
explores the experience of stress in university life and the ways
students can develop healthy approaches to managing stress and
supporting their well-being.

\section{What you will learn}\label{what-you-will-learn}

By working through this book, you will explore:

\begin{itemize}
\tightlist
\item
  Stress, stressors, and resilience.
\item
  The domains of self-care and the importance of social support.
\item
  The physiological and psychological effects of stress.
\item
  The relationships between stress, sleep, physical activity, and
  nutrition.
\item
  Mental health, diversity, and social justice.
\item
  Time and resource management.
\item
  Cognitive restructuring and reframing.
\item
  Spirituality, nature, and creative expression.
\item
  Strategies for developing a personal stress management plan.
\end{itemize}

\section{How to use this book}\label{how-to-use-this-book}

Read each chapter alongside the assigned course readings and learning
activities. Use the reflection questions to connect concepts with your
own experiences and identify strategies that may support your
well-being.

This book is a learning resource and does not replace professional
mental health care when it is needed.

\bookmarksetup{startatroot}

\chapter{Wellness, Stress, and
Resilience}\label{wellness-stress-and-resilience}

\section{Introduction}\label{introduction}

University life can be an exciting opportunity for learning, personal
growth, and new experiences. It can also involve competing
responsibilities, academic expectations, financial pressures, changing
relationships, and uncertainty about the future.

NURS 180 introduces students to foundational knowledge about stress,
stressors, resilience, social support, and strategies for promoting
physical, mental, emotional, and spiritual well-being.

This chapter introduces wellness, stress, resilience, and the experience
of transitioning to university. It provides a foundation for exploring
practical approaches to managing stress throughout the course.

\section{Learning objectives}\label{learning-objectives}

By the end of this chapter, students should be able to:

\begin{enumerate}
\def\labelenumi{\arabic{enumi}.}
\tightlist
\item
  Explain the concepts of wellness, stress, stressors, and resilience.
\item
  Identify common stressors experienced by university students.
\item
  Describe how stress can be experienced differently by different
  individuals.
\item
  Reflect on personal responses to stressful situations.
\item
  Identify initial strategies and sources of support that may help them
  navigate student life.
\end{enumerate}

\section{1.1 Understanding wellness}\label{understanding-wellness}

Wellness can be understood as an ongoing process of recognizing personal
needs and developing practices that support well-being.

Well-being is not simply the absence of illness or difficulty. Students
may experience challenges while also finding meaning, maintaining
relationships, learning new skills, and developing strategies to care
for themselves.

In this course, wellness includes physical, mental, emotional, and
spiritual dimensions of health. These dimensions are interconnected and
can be influenced by personal circumstances, relationships, and the
wider social environment.

For example, a student experiencing academic pressure may notice changes
in sleep, feel worried about an upcoming assignment, and have less time
to connect with friends. Considering these experiences together can help
the student identify what support or changes may be useful.

\section{1.2 Understanding stress and
stressors}\label{understanding-stress-and-stressors}

\textbf{Stress} refers to the experience of responding to demands,
challenges, or situations that a person perceives as requiring attention
or adaptation.

A \textbf{stressor} is a situation, event, demand, or circumstance that
contributes to a stress response.

Stressors may be related to academic responsibilities, finances, family
expectations, relationships, health, work, or uncertainty about the
future.

Stress is not always experienced in the same way. A situation that feels
manageable to one student may feel overwhelming to another. Previous
experiences, available resources, social support, and the meaning
attached to a situation may influence how a person responds.

\subsection{Common student stressors}\label{common-student-stressors}

Examples include:

\begin{itemize}
\tightlist
\item
  Managing multiple assignments and deadlines.
\item
  Preparing for examinations.
\item
  Adjusting to university expectations and new learning environments.
\item
  Balancing employment, studies, family, and personal responsibilities.
\item
  Experiencing financial uncertainty.
\item
  Navigating relationships, belonging, and social expectations.
\item
  Managing uncertainty about academic progress or future career plans.
\end{itemize}

This list is not exhaustive. Students may experience several stressors
at once, and their circumstances may change throughout the academic
year.

\subsection{Reflection: Identifying
stressors}\label{reflection-identifying-stressors}

Consider the following questions:

\begin{enumerate}
\def\labelenumi{\arabic{enumi}.}
\tightlist
\item
  What situations have contributed to stress in your academic or
  personal life?
\item
  Which stressors are within your control, partly within your control,
  or outside your immediate control?
\item
  What resources, relationships, or skills have helped you respond to
  these situations?
\end{enumerate}

You do not need to share private experiences with others. Choose an
example that you feel comfortable reflecting on.

\section{1.3 Thriving or surviving?}\label{thriving-or-surviving}

The distinction between thriving and surviving provides a useful
starting point for reflecting on student well-being.

\textbf{Surviving} may describe periods when a person is concentrating
mainly on getting through immediate demands. A student may complete
assignments and attend classes while feeling that there is little time
or energy left for rest, relationships, or activities they value.

\textbf{Thriving} may describe a greater sense of engagement,
connection, purpose, and capacity to manage challenges. Thriving does
not mean feeling happy or stress-free at all times.

These descriptions are not fixed categories. A student may feel that
they are thriving in one area of life while struggling in another. A
difficult week does not necessarily mean that a person is failing, and
experiencing stress does not mean that well-being is impossible.

The purpose of this reflection is to recognize current needs and
identify realistic opportunities for improvement.

\subsection{Reflection activity: Where am I
now?}\label{reflection-activity-where-am-i-now}

Think about your current experience of university life.

\begin{itemize}
\tightlist
\item
  What is going well for you?
\item
  What feels particularly demanding?
\item
  Are you making time for sleep, nourishment, movement, relationships,
  and activities that matter to you?
\item
  What is one manageable change that might support your well-being this
  week?
\end{itemize}

Record your reflections privately if you prefer.

\section{1.4 Understanding resilience}\label{understanding-resilience}

\textbf{Resilience} refers to the capacity and process of adapting to
challenges, responding to difficulties, and drawing on available
resources.

Resilience should not be understood as having to cope with everything
alone. Asking for help, taking a break, adjusting expectations, and
seeking appropriate support can all be part of responding constructively
to difficult circumstances.

Resilience can involve personal skills, relationships, community
resources, and the conditions in which a person lives and studies. It is
therefore important to consider both individual strategies and the
availability of supportive environments.

For example, a student who is struggling with a demanding workload might
review their schedule, speak with an instructor, seek academic support,
and ask a trusted person for help. These actions may help the student
respond to the situation more effectively than trying to manage every
difficulty alone.

\subsection{Building resilience through
reflection}\label{building-resilience-through-reflection}

Consider the following questions:

\begin{itemize}
\tightlist
\item
  How have you responded to a difficult situation in the past?
\item
  What helped you manage that experience?
\item
  What did you learn about your needs, strengths, or limits?
\item
  What would you do differently if a similar situation occurred again?
\item
  Who or what could support you during a future challenge?
\end{itemize}

Resilience does not require a person to tolerate unsafe, unfair, or
persistently harmful circumstances. Sometimes, responding to a challenge
requires changing the situation or seeking support rather than simply
adapting to it.

\section{1.5 Developing an initial stress-management
approach}\label{developing-an-initial-stress-management-approach}

Managing stress begins with recognizing what is happening and
considering what kind of response may be appropriate.

A useful starting process is:

\begin{enumerate}
\def\labelenumi{\arabic{enumi}.}
\tightlist
\item
  \textbf{Notice:} Identify the situation contributing to stress.
\item
  \textbf{Reflect:} Consider how the situation is affecting you and what
  you need.
\item
  \textbf{Prioritize:} Identify what requires immediate attention and
  what can wait.
\item
  \textbf{Choose:} Select one realistic action, such as breaking a task
  into smaller steps or contacting a support person.
\item
  \textbf{Review:} Reflect on whether the action helped and whether
  another response is needed.
\end{enumerate}

This is a practical reflection framework for the course, not a guarantee
that every stressful situation can be resolved through individual
action.

Different situations call for different responses. Some may benefit from
planning and problem-solving, while others may require emotional
support, rest, professional assistance, or changes to external
circumstances.

\section{1.6 Academic stress and the transition to
university}\label{academic-stress-and-the-transition-to-university}

The transition to university may involve changes in learning
expectations, independence, time management, relationships, and personal
responsibilities.

Students may need to develop new approaches to studying, seek
clarification when expectations are unclear, and learn how to balance
academic work with other aspects of life.

Academic integrity is also part of navigating university successfully.
Understanding course expectations, acknowledging sources, managing
deadlines, and seeking guidance when uncertain can help students
approach academic responsibilities constructively.

Students are encouraged to identify support resources available through
their institution, including academic advising, learning support,
student wellness services, and trusted instructors or peers.

Seeking assistance is not evidence of personal failure. It can be a
practical way of responding to a challenge.

\section{Chapter summary}\label{chapter-summary}

In this chapter, we explored wellness, stress, stressors, resilience,
and the experience of university life.

Key points include:

\begin{itemize}
\tightlist
\item
  Wellness involves several interconnected dimensions of health and
  well-being.
\item
  Stressors may arise from academic, financial, social, personal, and
  other circumstances.
\item
  People may experience and respond to similar situations differently.
\item
  Thriving and surviving are useful concepts for reflecting on current
  needs, not permanent labels.
\item
  Resilience can involve personal skills, supportive relationships, and
  access to resources.
\item
  Recognizing stressors and choosing realistic next steps can provide a
  starting point for stress management.
\item
  Seeking help can be an important part of navigating challenges.
\end{itemize}

\section{Review questions}\label{review-questions}

\begin{enumerate}
\def\labelenumi{\arabic{enumi}.}
\tightlist
\item
  What is the difference between stress and a stressor?
\item
  Identify three stressors that university students may experience.
\item
  How might two students respond differently to the same stressful
  situation?
\item
  What is resilience, and why does it not mean coping alone?
\item
  How might a student recognize that they are surviving rather than
  thriving?
\item
  Identify one realistic action that could help a student respond to an
  academic stressor.
\item
  What forms of support might be available to a student experiencing
  persistent stress?
\end{enumerate}

\section{Further reading}\label{further-reading}

The NURS 180 syllabus identifies the following resources relevant to
this chapter:

\begin{itemize}
\tightlist
\item
  Malette, N. (2020). \emph{It's All Good (Unless It's Not): Mental
  Health Tips \& Strategies for Your Undergrad Years}. UBC Press.
\item
  Kaufman (2018). \emph{What Is Stress Resilience and Can It Be
  Learned?}
\item
  Richardson, King, Garrett, and Wrench (2012). \emph{Thriving or Just
  Surviving?}
\item
  Moghimi et al.~(2023). \emph{Mental Health Challenges, Treatment
  Experiences, and Care Needs of Post-Secondary Students: A
  Cross-Sectional Mixed-Methods Study.}
\end{itemize}

Consult the course syllabus and assigned reading links for the full
publication details and access information.

\bookmarksetup{startatroot}

\chapter{}\label{section}

\bookmarksetup{startatroot}

\chapter{}\label{section-1}

\bookmarksetup{startatroot}

\chapter{}\label{section-2}

\bookmarksetup{startatroot}

\chapter{}\label{section-3}

\bookmarksetup{startatroot}

\chapter{}\label{section-4}

\bookmarksetup{startatroot}

\chapter{}\label{section-5}

\bookmarksetup{startatroot}

\chapter{}\label{section-6}

\bookmarksetup{startatroot}

\chapter{}\label{section-7}

\bookmarksetup{startatroot}

\chapter{}\label{section-8}

\bookmarksetup{startatroot}

\chapter{}\label{section-9}

\bookmarksetup{startatroot}

\chapter{}\label{section-10}


\backmatter


\end{document}
